\section{Задание}

\textbf{Цель работы:} исследование влияния \textit{свойств канала} связи на \textit{качество передачи сигналов} при различных методах физического и логического кодирования, используемых в цифровых сетях передачи данных.

В процессе выполнения учебно-исследовательской работы необходимо:
\begin{itemize}
    \item выполнить анализ качества передачи исходного сообщения по \textbf{идеальному} и \textbf{реальному} каналу связи при использовании разных методов кодирования и определить требуемую полосу пропускания канала, в том числе, для качественной передачи сообщения по реальному каналу в зависимости:
    \begin{enumerate}
        \item \textit{от уровня шумов и помех} в реальном канале;
        \item \textit{от степени рассинхронизации} передатчика и приёмника;
        \item \textit{уровня граничного напряжения}, в пределах которого невозможно распознать передаваемый сигнал из-за его затухания;
    \end{enumerate}
    \item выполнить \textbf{сравнительный анализ} рассматриваемых методов кодирования при передаче сообщения по идеальному и реальному каналу связи;
    \item выбрать и обосновать наилучший метод для передачи исходного сообщения по реальному каналу связи с учетом затухания, шумов в канале и рассинхронизации.
\end{itemize}

Исходный код к текущему отчёту можно найти в списке литературы.\cite{mrdvd2026itmo_labs}